% Pursuing Stacks (cotes 134-2 to 134-8): archive cote 134-7, site page 9 (PDF page 10), typescript page 451.
% typescript. Draft compared: ps6.tex:383-449.
% AI-transcribed by Claude Opus 5.5 (claude-opus-5-5, adaptive thinking, effort high), Message Batches API, batch ps3 (msgbatch_01XdFmuVVBAcgc7EFxdtCqsj), 2026-11-12; not yet checked by eye.
\scanpage{134-7/9}
\mspage{451}

\xout{xx})
(o) \qquad $L_k(X) = \deleted{\illegible{1}}\,\hand{\Gamma}_k^{\wedge}(M)\deleted{)}\,\inserted{\overset{\text{def}}{=}\prod_{i\geq 0}\Gamma^{i}(M)\deleted{)}}$ \hand{(}(k-linearization of X) ,
where M is any k-module endowed with an isomorphism
\[ u:\ X \simeq W(M) \quad , \]
the \inserted{k-module} \xout{object} $L_k(X)$ \uline{does not depend}, \uline{up to unique isomorphism}, \uline{on the
choice of a pair} (M,u), because \xout{anyxx} for two k-modules M,M', any
morphism of \uline{sheaves of sets}
\[ v:\ W(M) \longrightarrow W(M') \]
induces a\xout{n} \xout{isomorphism} homomorphism of k-modules
\[ \Gamma_k^{\wedge}(v):\quad \Gamma_k^{\wedge}(M) \longrightarrow \Gamma_k^{\wedge}(M') \]
(compatible not with multiplications, but with diagonal maps ...), \xout{whi}
which will be an isomorphism if v is.

Thus we do have, it seems to me, a good candidate for k-lineariza-%
tion. To check it is suitable indeed, the main point seems to check
that the corresponding diagram \hand{(l)} commutes up to canonical isomorphism,
the crucial case of course being $k=\underline{Z}$ . This now looks like a rather down
to earth question, which seems to me a pretty good test, whether the
intuition of schematization of homotopy types is a sound one. Let's
rephrase it here\over{,}{.} \over{f}{F}or this, let's first restate the \xout{definitionxx}
description of the \xout{\illegible{1}} category U(k) (coarse version) in the more
\xout{\illegible{1}}
down-to-earth terms of linear algebra. Objects may be viewed as just
k-modules M, whereas \xout{maps} (non additive\hand{)} ``maps''from M to M' (defined
previously as maps $W(M)\longrightarrow W(M')$ of sheaves of sets) are described as
just continuous k-linear maps
\[ \hand{(p)} \qquad \hand{f}:\ \Gamma_k^{\wedge}(M) \rightarrow \Gamma_k^{\wedge}(M') \quad , \]
which are mor\unsure{\hand{e}}over compatible with the natural augmentations to k,
and with the natural diagonal maps \unsure{\hand{:}|;}
\[ \hand{\varepsilon:\Gamma_k^{\wedge}(M)\rightarrow k,}\quad \hand{\Delta}:\ \Gamma_k^{\wedge}(M)\longrightarrow \Gamma_k^{\wedge}(M)\ \hand{\hat{\ }}\!\otimes_k\ \Gamma_k^{\wedge}(M) \quad (\hand{\simeq}\ \Gamma_k^{\wedge}(MxM)\ )\hand{,} \]
\inserted{(the latter,} \deleted{(}deduced from the usual linear diagonal map $M\rightarrow MxM$). When M is
looked \inserted{at} as being embedded in $\Gamma_k^{\wedge}(M)$ by the exponenti\xout{t}al map \xout{(xx}\hand{(}n),
it is identified (if I remember it right) to the set of elements \inserted{$\xi$} in
$\Gamma_k^{\wedge}(M)$ satisfying the relations
\[ \text{(q)} \qquad \hand{\varepsilon}(\hand{\xi}) = 1 \ ,\ \hand{\Delta}(\hand{\xi}) = \hand{\xi}\otimes\hand{\xi} \quad , \]
where \hand{$\varepsilon$} is the augmentation and \hand{$\Delta$} the diagonal map, hence p) induces
a map (in general not additive)

\margin{9}
